\documentclass[11pt]{article}
\usepackage[margin=1in]{geometry}

% Core packages
\usepackage{amsmath,amssymb}
\usepackage{array}
\usepackage{booktabs}
\usepackage{enumitem}
\usepackage{hyperref}
\usepackage{xcolor}

% Paragraphs
\setlength{\parindent}{0pt}
\setlength{\parskip}{0.8\baselineskip}
\emergencystretch=2em

% Lists
\setlist[itemize]{topsep=0.25\baselineskip,itemsep=0.25\baselineskip}
\setlist[enumerate]{topsep=0.25\baselineskip,itemsep=0.25\baselineskip}

% Table helpers
\newcolumntype{L}[1]{>{\raggedright\arraybackslash}p{#1}}

\title{Architecture Decision Record: ADR-001\\
\large Implementing OpenTelemetry Distributed Tracing}
\author{Abel-Raul Mazilu}
\date{13-05-2026}

\begin{document}
\maketitle

\begin{center}
\begin{tabular}{rl}
\textbf{Status:} & Accepted \\
\textbf{Project:} & DVerse Communication Platform 2.0 \\
\end{tabular}
\end{center}

\section{Context}

The DVerse platform is evolving toward a decoupled multi-agent architecture in which AI agents operate as autonomous participants. These agents communicate asynchronously over Zenoh and also interact with users through the chat application interface.

As this architecture scales, the absence of distributed tracing becomes an engineering liability. Without end-to-end trace visibility, it is difficult to understand how a user request, agent action, Zenoh publication, message consumption, and downstream agent response are connected. Failures become difficult to isolate, latency sources remain opaque, and behavior across agent boundaries is not reproducible.

Sprint 4 must therefore introduce tracing at the same time as agent communication infrastructure. Retrofitting tracing later would likely produce incomplete coverage and inconsistent instrumentation.

\section{Decision}

Sprint 4 will implement OpenTelemetry as the distributed tracing foundation for the multi-agent system. Traces will capture agent lifecycle events, Zenoh message publication and consumption, chat interface interactions, and cross-agent workflows.

The implementation will establish trace spans for the following events:

\begin{itemize}
    \item Agent startup, initialization, shutdown, and error states.
    \item Message publication to Zenoh topics.
    \item Message consumption from Zenoh subscriptions.
    \item Agent-to-agent request and response flows.
    \item Agent interactions that enter or leave the chat application.
\end{itemize}

Trace context must be propagated across Zenoh message boundaries. Because Zenoh does not provide HTTP-style headers for OpenTelemetry context propagation, the team will define a message-level propagation convention. This convention may include embedding trace context inside structured message metadata so that downstream agents can continue the same distributed trace.

\section{Alternative Decisions}

One alternative was to defer tracing to a later sprint. This was rejected because distributed tracing is most reliable when introduced alongside the communication paths it observes.

Another alternative was to use a vendor-specific tracing solution. This was rejected because OpenTelemetry provides an open standard, reduces vendor lock-in, and aligns with the broader cloud-native ecosystem.

\section{Consequences}

\subsection{Positive}

\begin{itemize}
    \item Engineers will be able to follow a single interaction across agent boundaries, Zenoh topics, and chat application integration points.
    \item Latency analysis and error diagnosis will become easier as the number of agents and topics grows.
    \item Future observability tools can consume OpenTelemetry data without requiring a redesign of instrumentation points.
\end{itemize}

\subsection{Negative}

\begin{itemize}
    \item Instrumentation adds implementation complexity during Sprint 4.
    \item Trace context propagation across Zenoh messages requires a documented convention.
    \item Overly granular tracing can introduce performance overhead in high-throughput message paths.
\end{itemize}

\section{Scope}

\begin{table}[ht]
\centering
\renewcommand{\arraystretch}{1.25}
\begin{tabular}{L{0.32\textwidth} L{0.58\textwidth}}
\toprule
\textbf{Deliverable} & \textbf{Description} \\
\midrule
Agent lifecycle tracing & Instrument startup, initialization, shutdown, and failure events. \\
Zenoh tracing & Trace publish and subscribe operations across Zenoh topics. \\
Chat tracing & Trace agent interactions entering or leaving the chat interface. \\
Trace propagation & Define and document message-level trace context propagation. \\
\bottomrule
\end{tabular}
\end{table}


\newpage

\section*{Statement on AI Use}

Artificial intelligence was used to support the drafting of parts of this document. The generated text was subsequently reviewed, checked, and edited where necessary by me, the author, Abel-Raul Mazilu.

\end{document}